\documentclass[12pt]{article}
\usepackage[spanish]{babel}
\usepackage[utf8]{inputenc}
\usepackage[T1]{fontenc}
\usepackage{amsmath, amssymb}
\usepackage{geometry}
\usepackage{enumitem}
\usepackage{needspace}
\usepackage{tikz}
\usetikzlibrary{babel}

\geometry{letterpaper, margin=2cm}
\setlength{\parindent}{0pt}
\setlength{\parskip}{6pt}

\newcommand{\puntografico}{0.09}

\begin{document}

\begin{center}
  {\Large \textbf{Borrador de items}}\\
  {\large Bloque 2. Relaciones y Algebra - Afirmacion 3}\\[4pt]
  Prueba Nacional Estandarizada de Matematica, Secundaria 2026
\end{center}

\section*{Afirmacion}
Resuelve problemas relacionados con la funcion dada por $f(x)=a\sqrt{x+b}+c$, en su representacion grafica o algebraica.

\section*{Items propuestos}

\begin{enumerate}[label=\textbf{\arabic*.}, itemsep=1.05cm]

\Needspace{0.58\textheight}
\item
Considere la siguiente informacion:

En un centro de acopio se estima la cantidad maxima $M$ de cajas pequenas que se pueden colocar en una zona de almacenamiento. Esa cantidad esta dada por
\[
M(x)=2\sqrt{x+9}+6,
\]
donde $x$ representa el area disponible adicional, en metros cuadrados, con $7 \leq x \leq 72$.

De acuerdo con la informacion anterior, \textquestiondown cual es la cantidad maxima de cajas pequenas que se pueden colocar si el area disponible adicional es de $16\text{ m}^2$?

\begin{enumerate}[label=\Alph*)]
  \item 11
  \item 14
  \item 16
  \item 31
\end{enumerate}

\Needspace{0.60\textheight}
\item
Considere la siguiente informacion:

En una planta de tratamiento se usa un sensor para estimar la cantidad $L$ de litros de agua filtrada por minuto. Esa cantidad esta dada por
\[
L(t)=3\sqrt{t+4}+2,
\]
donde $t$ representa el tiempo de funcionamiento del filtro, en minutos, con $0 \leq t \leq 60$.

De acuerdo con la informacion anterior, \textquestiondown cuantos minutos debe funcionar el filtro para que la cantidad estimada sea de $17$ litros por minuto?

\begin{enumerate}[label=\Alph*)]
  \item 9
  \item 21
  \item 25
  \item 29
\end{enumerate}

\Needspace{0.96\textheight}
\item
Considere la siguiente informacion:

En un laboratorio escolar se registra la altura $h(x)$, en centimetros, de una columna de espuma que se forma durante una reaccion. La variable $x$ representa el tiempo, en minutos, desde que inicia la observacion.

La siguiente representacion grafica muestra el comportamiento de esa altura:

\begin{center}
\begin{tikzpicture}[xscale=0.58, yscale=0.55]
  \path[use as bounding box] (-0.9,-0.8) rectangle (11.2,10.8);
  \draw[->, thick] (-0.5,0) -- (10.6,0) node[anchor=west] {$x$};
  \draw[->, thick] (0,-0.4) -- (0,10.2) node[anchor=south] {$h(x)$};

  \foreach \x/\lab in {1/1,2/2,5/5,10/10}
    \node[anchor=north] at (\x,0) {$\lab$};
  \foreach \y/\lab in {3/3,5/5,7/7,9/9}
    \node[anchor=east] at (0,\y) {$\lab$};

  \draw[densely dashed, gray!80] (1,0) -- (1,3) -- (0,3);
  \draw[densely dashed, gray!80] (2,0) -- (2,5) -- (0,5);
  \draw[densely dashed, gray!80] (5,0) -- (5,7) -- (0,7);
  \draw[densely dashed, gray!80] (10,0) -- (10,9) -- (0,9);

  \draw[thick, domain=1:10, samples=120, smooth, variable=\x]
    plot ({\x},{2*sqrt(\x-1)+3});

  \fill (1,3) circle (\puntografico);
  \fill (2,5) circle (\puntografico);
  \fill (5,7) circle (\puntografico);
  \fill (10,9) circle (\puntografico);
\end{tikzpicture}
\end{center}

De acuerdo con la informacion anterior, \textquestiondown cual representacion algebraica corresponde a la funcion $h$?

\begin{enumerate}[label=\Alph*)]
  \item $h(x)=2\sqrt{x-1}+3$
  \item $h(x)=2\sqrt{x+1}+3$
  \item $h(x)=2\sqrt{x-1}-3$
  \item $h(x)=\sqrt{x-1}+3$
\end{enumerate}

\end{enumerate}

\section*{Clave y justificacion breve}

\begin{enumerate}[label=\textbf{\arabic*.}]
  \item \textbf{C.} Al sustituir $x=16$, se obtiene $M(16)=2\sqrt{16+9}+6=2\sqrt{25}+6=16$.
  \item \textbf{B.} Si $L(t)=17$, entonces $3\sqrt{t+4}+2=17$, por lo que $\sqrt{t+4}=5$ y $t+4=25$. Por tanto, $t=21$.
  \item \textbf{A.} La grafica inicia en $(1,3)$ y pasa por puntos como $(2,5)$ y $(5,7)$; esto corresponde a $h(x)=2\sqrt{x-1}+3$.
\end{enumerate}

\end{document}
